\documentclass[12pt]{article}
\usepackage[utf8]{inputenc}
\usepackage[T1]{fontenc}
\usepackage{hyperref}
\title{Puluwatese Language Resources for K-12 Curriculum Development}
\author{Compiled by Claude Sonnet 5.5 (claude-sonnet-5-5) for Andr\'as Kornai}
\date{}
\begin{document}
\maketitle

\section{Introduction}
Puluwatese (ISO 639-3 \texttt{puw}) is a Micronesian (Austronesian) language of the Federated States of Micronesia, spoken on Puluwat atoll. Counts are about 1,400 (1989 census) to 1,500 (2002, per Omniglot). Dialects listed include Pulapese and Pulusukese. UNESCO classifies it as severely endangered per the Wikipedia article. Only linguistic reference material was found.

\section{Primary Educational Resources}
None found in these searches.

\section{Linguistic and Reference Materials}
\begin{itemize}
\item Elbert, \emph{Puluwat Dictionary} (Pacific Linguistics C-24, 1972), Puluwatese--English, freely available (CC BY-SA 4.0 online edition): \url{https://openresearch-repository.anu.edu.au/bitstream/1885/146583/1/PL-C24.pdf}
\item Library catalogue record for the dictionary: \url{https://find.library.upenn.edu/catalog/994760593503681}
\item Wikipedia, Puluwat language: \url{https://en.wikipedia.org/wiki/Puluwat_language}
\item Omniglot page on Puluwatese: \url{https://omniglot.com/writing/puluwatese.htm}
\item Austronesian Comparative Dictionary source entry for Elbert 1972: \url{https://acd.clld.org/sources/elbert1972}
\end{itemize}

\section{Religious and Community Resources}
None found in these searches.

\section{Digital and Audio Resources}
\begin{itemize}
\item Verbix documentation page on Puluwatese: \url{https://docs.verbix.com/Languages/Puluwatese}
\end{itemize}

\section{Curriculum Development Context}
Puluwatese is spoken by all in its community, but no school materials were found. Elbert's dictionary is the key resource for vocabulary and orthographic conventions of its time (1972); current practice should be checked with the community.

\end{document}